\section{Diagnósticos matemáticos independientes de la búsqueda de ocultedad}
\label{sec:diagnosticos-independientes}

Las funciones de esta sección aceptan una trayectoria o una serie temporal
suministrada por la persona usuaria. No ejecutan por sí mismas una búsqueda de
atractores ocultos y sus salidas no contienen una decisión de cuenca. Para una
trayectoria uniformemente muestreada se usa
\[
 t_n=t_0+n\Delta t,\qquad
 X_n=(x_{n,1},\ldots,x_{n,d})^\mathsf{T},\qquad n=0,\ldots,N-1.
\]
Cuando se declara un tiempo de descarte \(t_{\rm burn}\), todas las métricas se
calculan sobre el conjunto de índices
\(\mathcal I=\{n:t_n\geq t_{\rm burn}\}\). Las etiquetas devueltas son
diagnósticos finito-temporales; no son pruebas asintóticas de acotamiento,
caos, periodicidad u ocultedad.

\subsection{Métricas de trayectoria y acotamiento}

Para cada componente, \codepath{compute_trajectory_metrics} calcula el rango y
la varianza de la cola:
\[
 R_j=\max_{n\in\mathcal I}x_{n,j}
     -\min_{n\in\mathcal I}x_{n,j},
 \qquad
 s_j^2=\frac{1}{|\mathcal I|}
 \sum_{n\in\mathcal I}(x_{n,j}-\bar x_j)^2.
\]
Si se proporcionan equilibrios \(E_\ell\), también calcula
\[
 d_{\min}(X_{N-1})=\min_\ell\|X_{N-1}-E_\ell\|_2
\]
y sólo registra contacto final cuando esa distancia no excede
\codepath{equilibrium_tol}. La etiqueta \codepath{bounded_nontrivial} significa
únicamente que la trayectoria es finita, no rebasó el radio numérico de
divergencia y no terminó dentro de la tolerancia de un equilibrio.

\codepath{compute_boundedness_metrics} forma la serie
\[
 r_n=
 \begin{cases}
 \|X_n\|_2,&\texttt{norm="euclidean"},\\
 \|X_n\|_\infty,&\texttt{norm="max"},
 \end{cases}
\]
y, con una fracción de ventana \(\gamma\in(0,0.5]\), usa
\[
 G=\frac{\operatorname{mean}(r_{N-m},\ldots,r_{N-1})}
 {\max\{\operatorname{mean}(r_0,\ldots,r_{m-1}),\epsilon_{\rm mach}\}},
 \qquad m=\max\{2,\operatorname{round}(\gamma N)\}.
\]
La implementación clasifica \codepath{unbounded_candidate} si se rebasa
\codepath{divergence_radius} o si \(G\geq100\); clasifica
\codepath{bounded_candidate} si el máximo es finito y \(G<10\); los valores
intermedios son inconclusos. Estos umbrales son reglas operativas de la
librería, no teoremas de estabilidad.

\begin{verbatim}
from hidden_attractors import (
    compute_boundedness_metrics,
    compute_trajectory_metrics,
    trajectory_metrics_for_system,
)

trajectory_metrics = compute_trajectory_metrics(
    times, states, equilibria=equilibria, t_start=40.0
)
system_metrics = trajectory_metrics_for_system(
    trajectory, system=system, h=0.01, t_start=40.0
)
boundedness = compute_boundedness_metrics(
    times, states, burn_time=40.0,
    divergence_radius=120.0,
)
\end{verbatim}

\codepath{trajectory_metrics_for_system} es una fachada para las mismas
ecuaciones: acepta una matriz \((t,X)\) y obtiene los equilibrios del objeto
\codepath{system}, o acepta el diccionario \codepath{equilibria} de forma
explícita. Si \codepath{has_time=False}, construye
\(t_n=nh\); si la matriz ya contiene tiempo, \(h\) no modifica sus muestras.
La función exige una de las dos fuentes de equilibrios y no ejecuta búsqueda
de semillas ni una prueba de ocultedad.

\begin{verbatim}
trajectory_metrics_for_system(
    traj, *, system=None, equilibria=None, h, t_start,
    divergence_norm=120.0, equilibrium_tol=1.0e-3,
    has_time=True,
)
\end{verbatim}

\subsection{DFT, FFT, potencia, entropía espectral y Welch}

Después de retirar tendencia o media y aplicar una ventana \(w_n\), la
transformada discreta implementada por NumPy es
\[
 Y_k=\sum_{n=0}^{N-1}w_n y_n
       \exp\!\left(-2\pi\ii\frac{kn}{N}\right),
 \qquad
 f_k=\frac{k}{N\Delta t},
 \qquad 0\leq k\leq\left\lfloor\frac N2\right\rfloor .
\]
\codepath{fft_spectrum} devuelve la amplitud unilateral
\[
 A_k=\frac{|Y_k|}{\sum_n w_n},
\]
mientras que \codepath{compute_fft_psd} usa potencia no normalizada
\(P_k=|Y_k|^2\) y, por defecto, la normaliza como
\(p_k=P_k/\sum_jP_j\). De esa distribución obtiene
\[
 H_{\rm spec}=
 -\frac{\sum_{k:p_k>0}p_k\log p_k}{\log K},
 \qquad
 D_{\rm peak}=\frac{\max_k p_k}{\sum_kp_k}.
\]
La fracción de ancho de banda es el número mínimo de bins de mayor potencia
que acumulan el \(90\) por ciento, dividido por el número total de bins. La
clasificación de pico dominante, candidato cuasiperiódico o espectro de banda
ancha usa umbrales explícitos del código y no certifica caos. La FFT rápida es
el algoritmo computacional para evaluar la DFT
\cite{CooleyTukey1965}; la entropía normalizada deriva de la entropía de
Shannon \cite{Shannon1948}.

Para \codepath{psd_welch}, cada segmento solapado \(r\), de longitud \(L\), se
centra y multiplica por una ventana de Hann. La estimación implementada es
\[
 \widehat S_{yy}(f_k)=
 \frac1{M}\sum_{r=1}^{M}
 c_k\frac{\Delta t\left|\sum_{n=0}^{L-1}
 w_n y_{r,n}\exp(-2\pi\ii kn/L)\right|^2}
 {\sum_{n=0}^{L-1}w_n^2},
\]
donde \(c_k=2\) en los bins interiores positivos y \(c_k=1\) en continua y,
para longitud par, Nyquist. Esta es la densidad unilateral de Welch
\cite{Welch1967}; si \(y\) tiene unidades \(U\) y el tiempo unidades \(T\), la
salida tiene unidades \(U^2/\mathrm{Hz}=U^2T\). La implementación se contrasta
con \codepath{scipy.signal.welch} usando la misma ventana explícita, solapamiento
y detrend. La entrada debe estar uniformemente muestreada; aumentar el
solapamiento no crea información y la ventana intercambia fuga espectral por
resolución.

\begin{verbatim}
from hidden_attractors import compute_fft_psd
from hidden_attractors.analysis import fft_spectrum, psd_welch

summary = compute_fft_psd(
    times, states[:, 0], burn_time=40.0,
    window="hann", normalize_power=True,
)
fft = fft_spectrum(states[:, 0], h=times[1] - times[0])
welch = psd_welch(
    states[:, 0], h=times[1] - times[0],
    nperseg=512, overlap=0.5,
)
\end{verbatim}

\subsection{Cruces de una sección de Poincaré}

\par\smallskip\noindent
Para la sección \(x_s=c\), la función
\codepath{detect_poincare_crossings} acepta un cruce positivo cuando
\[
 x_{n,s}<c\leq x_{n+1,s},
 \qquad x_{n+1,s}-x_{n,s}>0,
\]
y un cruce negativo cuando
\[
 x_{n,s}>c\geq x_{n+1,s},
 \qquad x_{n+1,s}-x_{n,s}<0.
\]
Estas desigualdades evitan contar dos veces una muestra situada exactamente
en la sección. Con interpolación lineal calcula
\[
 \theta_n=\frac{c-x_{n,s}}{x_{n+1,s}-x_{n,s}},\qquad
 t_\star=t_n+\theta_n(t_{n+1}-t_n),\qquad
 X_\star=X_n+\theta_n(X_{n+1}-X_n).
\]
En \codepath{derivative_mode="integer_rhs"}, además se exige que
\(F_s(t_\star,X_\star)\) tenga el signo de cruce. En
\codepath{derivative_mode="geometric_fractional"} sólo se usa la orientación
de las muestras. Por tanto, para Caputo la salida es una sección geométrica
numérica y no un mapa de retorno clásico exacto. Tampoco el modo entero es un
mapa exacto: sigue siendo una interpolación lineal entre muestras; la
evaluación de \(F\) únicamente aporta la prueba clásica de orientación. En
consecuencia, la API registra \codepath{exact_poincare_map=False} en todos los
modos y distingue el caso entero mediante
\codepath{classical_integer_section_interpretation=True}. La separación entre
localización del evento y orientación sigue la práctica de cálculo numérico de
mapas de Poincaré \cite{Henon1982}.

\begin{verbatim}
from hidden_attractors import detect_poincare_crossings

section = detect_poincare_crossings(
    times,
    states,
    section_variable="x",
    section_value=0.0,
    direction="positive",
    derivative_mode="geometric_fractional",
    burn_time=40.0,
)
\end{verbatim}

\subsection{Prueba 0--1 de Gottwald--Melbourne}

Para una señal escalar \(\phi_j\) y un valor
\(c\in(\pi/5,4\pi/5)\), la implementación construye
\[
 p_c(n)=\sum_{j=1}^{n}\phi_j\cos(jc),\qquad
 q_c(n)=\sum_{j=1}^{n}\phi_j\sin(jc).
\]
Para un retardo \(\ell\), calcula el desplazamiento cuadrático medio
\[
 M_c(\ell)=\frac1{N-\ell}
 \sum_{j=1}^{N-\ell}
 \left([p_c(j+\ell)-p_c(j)]^2+
       [q_c(j+\ell)-q_c(j)]^2\right),
\]
retira el término oscilatorio
\[
 V_{\rm osc}(\ell)=
 \bar\phi^2\frac{1-\cos(\ell c)}{1-\cos c},
 \qquad
 D_c(\ell)=M_c(\ell)-V_{\rm osc}(\ell),
\]
y usa la correlación
\[
 K_c=\operatorname{corr}\bigl((1,\ldots,L),
 (D_c(1),\ldots,D_c(L))\bigr).
\]
El valor reportado es la mediana de \(K_c\) sobre valores reproducibles de
\(c\), siguiendo la formulación robusta de
Gottwald y Melbourne \cite{GottwaldMelbourne2004,GottwaldMelbourne2009}.
La librería usa \(K>0.8\) como candidato caótico, \(K<0.2\) como candidato
regular y la zona intermedia como inconclusa. Ruido, tendencia, sobremuestreo
y transitorios pueden sesgar el resultado.

\begin{verbatim}
from hidden_attractors import zero_one_test

result_01 = zero_one_test(
    states[:, 0],
    n_c=100,
    random_seed=12345,
    detrend=True,
    normalize=True,
)
\end{verbatim}

\subsection{Postprocesamiento para diagramas de bifurcación}

La función \codepath{bifurcation_points_from_trajectories} no realiza
continuación de ramas. Para cada parámetro \(\mu_r\), descarta el transitorio y
selecciona máximos locales de un observable \(y_n\):
\[
 y_n-y_{n-1}\geq0,\qquad y_n-y_{n+1}\geq0.
\]
Los modos \codepath{minima}, \codepath{both} y \codepath{sample} cambian esa
regla. El diagrama es el conjunto
\(\{(\mu_r,y_{r,n}):n\in\mathcal M_r\}\); no identifica automáticamente
puntos de bifurcación, estabilidad de ramas ni multiplicadores de Floquet.
Para un conjunto no vacío \(\mathcal B\) de puntos extraídos,
\codepath{bifurcation_summary} sólo resume
\[
 n_{\mathcal B}=|\mathcal B|,\qquad
 \mu_{\min}=\min_{b\in\mathcal B}\mu_b,\quad
 \mu_{\max}=\max_{b\in\mathcal B}\mu_b,
\]
\[
 y_{\min}=\min_{b\in\mathcal B}y_b,\qquad
 y_{\max}=\max_{b\in\mathcal B}y_b.
\]
Si \(\mathcal B\) está vacío devuelve únicamente
\codepath{n_points=0}; este resumen tampoco localiza bifurcaciones.
Su firma pública es \codepath{bifurcation_summary(points)}.

\begin{verbatim}
from hidden_attractors import (
    bifurcation_points_from_trajectories,
    bifurcation_summary,
)
from hidden_attractors.plotting import plot_bifurcation_diagram

points = bifurcation_points_from_trajectories(
    scans, observable="x", t_start=40.0, mode="maxima"
)
point_summary = bifurcation_summary(points)
plot_bifurcation_diagram(
    points, "bifurcation.pdf",
    parameter_label="mu",
    observable_label="maximos locales de x",
)
\end{verbatim}

\subsection{Exponentes de Lyapunov desde las ecuaciones}

\subsubsection{QR--Benettin para \(q=1\)}

Para \(\dot X=F(X)\), la matriz variacional satisface
\[
 \dot\Phi(t)=J(X(t))\Phi(t),\qquad
 J(X)=\frac{\partial F}{\partial X},\qquad \Phi(0)=I.
\]
Después de cada intervalo de reortogonalización,
\(\Phi_k=Q_kR_k\), y el estimador finito-temporal es
\[
 \widehat\lambda_i(T)=
 \frac1T\sum_{k=1}^{m}\log |(R_k)_{ii}|.
\]
Ésta es la ruta \codepath{integer_qr_benettin}, basada en el método de
Benettin y la práctica de reortogonalización para espectros numéricos
\cite{Benettin1980a,Benettin1980b,Wolf1985}. Sólo es válida para \(q=1\).
Si no se proporciona Jacobiano analítico, la librería usa diferencias finitas
y registra la advertencia correspondiente. La discretización implementada
debe quedar explícita:
\[
 X_{n+1}=\operatorname{EFORK}_{q=1}(F,X_n,h),\qquad
 \Phi_{n+1}=\Phi_n+hJ(X_n)\Phi_n.
\]
Es decir, el estado usa el paso de tres etapas
\codepath{efork_q1_step}, mientras que la base tangente usa Euler explícito
antes de cada QR. La exactitud de los exponentes debe comprobarse refinando
\(h\); la etiqueta QR--Benettin identifica la construcción variacional y la
reortogonalización, no un mismo integrador de alto orden para ambos bloques.
\codepath{integer_qr_benettin_lyapunov_exponents} es la fachada canónica
directa para \(F\), \(J\) y \(X_0\); reutiliza el mismo núcleo y no constituye
un segundo método. Impone el contrato entero
\(\lvert q-1\rvert\leq10^{-9}\).
\codepath{integer_system_lyapunov_exponents} adapta un
\codepath{ChaoticSystem}: usa su \codepath{evaluate}, emplea
\codepath{jacobian_matrix} cuando el sistema declara un Jacobiano analítico y,
en caso contrario, aproxima
\[
 J_{ij}(X)\approx
 \frac{F_i(X+\varepsilon e_j)-F_i(X-\varepsilon e_j)}
 {2\varepsilon}.
\]
La fachada intenta inferir el orden declarado del sistema y lo rechaza cuando
\(\lvert q_{\mathrm{system}}-1\rvert>10^{-9}\). Si el objeto no declara orden, lo acepta por
compatibilidad; esa aceptación no demuestra que el modelo sea entero.

\begin{verbatim}
integer_qr_benettin_lyapunov_exponents(
    rhs, jacobian, x0, *, h, t_final, t_burn=0.0,
    reorthonormalize_every=10, jacobian_eps=1.0e-6,
    div_threshold=None, q=1.0,
)
integer_system_lyapunov_exponents(
    system, x0, *, h, t_final, t_burn=0.0,
    reorthonormalize_every=10, jacobian_eps=1.0e-6,
    div_threshold=None,
)
\end{verbatim}

\begin{verbatim}
from hidden_attractors import (
    integer_qr_benettin_lyapunov_exponents,
    integer_system_lyapunov_exponents,
)

direct = integer_qr_benettin_lyapunov_exponents(
    rhs, jacobian, x0, h=0.01, t_final=20.0,
    t_burn=2.0, q=1.0,
)
registered = integer_system_lyapunov_exponents(
    system, x0, h=0.01, t_final=20.0, t_burn=2.0,
)
\end{verbatim}

\subsubsection{Sistema variacional Caputo ABM--QR}

Para \(0<q<1\), la ruta
\codepath{fractional_variational_abm_qr} integra el sistema extendido
\[
 \CaputoD X=F(X),\qquad
 \CaputoD\Phi=J(X)\Phi,\qquad \Phi(0)=I
\]
con ABM. Como el operador de Caputo depende de la historia, después de
\(\Phi_k=Q_kR_k\) la implementación coherente transforma cada bloque
variacional almacenado:
\[
 \Phi_j\leftarrow \Phi_jR_k^{-1},
 \qquad
 G_j\leftarrow
 \begin{bmatrix}F(X_j)\\J(X_j)\Phi_j\end{bmatrix},
 \quad j\leq k.
\]
Con \codepath{memory_mode="full"} se conserva toda la historia; una ventana
finita cambia el problema numérico. El sistema original--variacional se apoya
en la formulación discutida por Danca y Kuznetsov
\cite{DancaKuznetsov2018}; la transformación de \emph{toda} la historia
almacenada después de QR es una extensión propia de esta librería y no una
reproducción literal del algoritmo publicado. La implementación tiene pruebas
sintéticas, pero no se presenta como validación cuantitativa publicada para
todos los modelos, órdenes o superficies no suaves.

\subsubsection{Dinámica clonada}

Las dos rutas clonadas no requieren Jacobiano. En cada bloque se integran una
trayectoria fiducial y \(d\) copias
\[
 X^{(i)}_0=X_0+\delta e_i,\qquad i=1,\ldots,d.
\]
Al final del bloque se forma
\[
 V_k=\left[
 X^{(1)}_k-X^{(0)}_k,\ldots,
 X^{(d)}_k-X^{(0)}_k
 \right],
\]
se ortogonaliza mediante Gram--Schmidt o QR, y se acumulan las normas
\(\rho_{k,i}\):
\[
 \widehat\lambda_i=
 \frac1{K\,t_{\rm clone}}
 \sum_{k=1}^{K}\log\frac{\rho_{k,i}}{\delta}.
\]
\codepath{fractional_cloned_dynamics_abm_gs_published} sigue la ruta de
Fischer, Zourmba y Mohamadou \cite{Fischer2020}; la variante
\codepath{fractional_cloned_dynamics_abm_qr} sustituye GS por QR. Ambas usan
la misma semántica de reinicio ABM por bloques para la integración
fraccionaria. La ejecución registra
\codepath{effective_memory_protocol="published_block_restart"}; los nombres
anteriores, como \codepath{experimental_qr_block_restart}, sólo se conservan
en \codepath{requested_memory_protocol} como alias de compatibilidad. La
diferencia efectiva de la variante es la ortogonalización QR, no otra historia
de Caputo. Ninguna de las dos es un estimador Caputo de memoria completa.

\begin{center}
\footnotesize
\begin{longtable}{@{}L{0.27\textwidth}L{0.12\textwidth}L{0.13\textwidth}L{0.20\textwidth}L{0.20\textwidth}@{}}
\toprule
Método solicitado & Orden & Jacobiano & Memoria & Estado de evidencia \\
\midrule
\codepath{integer_qr_benettin} & \(q=1\) & Opcional; DF centradas & No aplica & Controles lineales exactos e intercomparación interna; sin reproducción cuantitativa de un espectro publicado.\\
\codepath{fractional_variational_abm_qr} & \(0<q<1\) & Opcional; DF centradas & Completa o ventana declarada & Implementado; validación sintética, sin claim cuantitativo publicado general.\\
\codepath{fractional_cloned_dynamics_abm_gs_published} & \(0<q\leq1\) & No & Reinicio por bloques & Reproducción con discrepancia de benchmark registrada.\\
\codepath{fractional_cloned_dynamics_abm_qr} & \(0<q\leq1\) & No & Reinicio por bloques & Variante numérica para comparación; sin claim publicado.\\
\bottomrule
\end{longtable}
\end{center}

\subsubsection{Contrato explícito de validación del método}

\codepath{validate_lyapunov_method_request} comprueba una
\codepath{LyapunovComputationRequest} sin ejecutar el integrador y devuelve
exactamente la terna
\[
 (\mathrm{ok},\mathrm{status},\mathrm{warnings}).
\]
Para una solicitud compatible devuelve
\codepath{(True, "compatible", warnings)}; una incompatibilidad ordinaria se
representa con \codepath{ok=False} y no mediante una excepción. Primero exige
una fuente y sólo una para el campo vectorial: exactamente una de
\codepath{system} o \codepath{rhs}:
\[
 \mathbf{1}_{\{\mathrm{system}\neq\mathrm{None}\}}+
 \mathbf{1}_{\{\mathrm{rhs}\neq\mathrm{None}\}}=1.
\]
También exige
\[
 X_0\in\mathbb{R}^{d},\ d\geq1,\quad
 h>0,\quad T_{\mathrm{final}}>0,\quad T_{\mathrm{burn}}\geq0,\quad
 \varepsilon_J>0,
\]
con \(X_0\), \(q\), los tiempos, \(h\) y \(\varepsilon_J\) finitos.
Si se proporciona, \codepath{reorthonormalize_every} debe ser un entero
positivo. \codepath{memory_window} sólo se inspecciona en la ruta variacional
con \codepath{memory_mode="window"} y entonces también debe ser un entero
positivo. Asimismo,
\codepath{reorthonormalization_time} y \codepath{div_threshold} deben ser
finitos y positivos.

La validación específica reproduce el registro de métodos: QR--Benettin
entero requiere \(q=1\) y memoria \codepath{not_applicable}; ABM--QR
variacional requiere \(0<q<1\), memoria \codepath{full} o
\codepath{window}, y comprueba la igualdad entre el \(q\) solicitado y el
declarado por el sistema cuando ambos existen. Las dos rutas de dinámica
clonada requieren \(0<q\leq1\) y aceptan
\codepath{not_applicable}, \codepath{published_block_restart} o el alias
\codepath{experimental_qr_block_restart}. La ausencia de Jacobiano en los
métodos variacionales produce
\codepath{analytic_jacobian_missing_finite_difference_used}; las rutas
clonadas producen \codepath{cloned_dynamics_no_jacobian_required}. Son
advertencias, no certificaciones de caos, ocultedad o validez fraccionaria.

Los estados de incompatibilidad que puede devolver el código actual son:
\begin{itemize}
  \item \codepath{unknown_method};
  \item \codepath{invalid_parameter};
  \item \codepath{method_not_valid_for_fractional_caputo};
  \item \codepath{memory_mode_not_applicable_for_integer_method};
  \item \codepath{method_not_valid_for_integer_or_out_of_range_q};
  \item \codepath{memory_mode_must_be_full_or_window_for_fractional_method};
  \item \codepath{request_q_does_not_match_system_q};
  \item \codepath{method_not_valid_for_out_of_range_q};
  \item \codepath{memory_mode_must_be_block_restart_for_cloned_dynamics}.
\end{itemize}

\begin{verbatim}
from hidden_attractors import (
    LyapunovComputationRequest,
    validate_lyapunov_method_request,
)

request = LyapunovComputationRequest(
    system=system, rhs=None, jacobian=None, x0=x0,
    q=1.0, method="integer_qr_benettin",
    h=0.01, t_final=20.0, t_burn=2.0,
    reorthonormalize_every=10,
    memory_mode="not_applicable",
)
ok, status, warnings = validate_lyapunov_method_request(request)
\end{verbatim}

\begin{verbatim}
import numpy as np
from hidden_attractors import compute_lyapunov_spectrum

rhs = lambda x: np.array([-x[0], -2.0*x[1]])
jac = lambda x: np.array([[-1.0, 0.0], [0.0, -2.0]])

summary = compute_lyapunov_spectrum(
    rhs=rhs,
    jacobian=jac,
    x0=np.array([1.0, 1.0]),
    q=1.0,
    method="integer_qr_benettin",
    h=0.01,
    t_final=20.0,
    t_burn=2.0,
)
print(summary.result.exponents)
\end{verbatim}

\subsection{Lyapunov desde una serie temporal escalar}

Esta ruta es distinta de los métodos variacionales. A partir de una señal
\(s_n\), una reconstrucción por retardos forma \cite{Takens1981}
\[
 Y_i=(s_i,s_{i+\tau},\ldots,s_{i+(m-1)\tau})^\mathsf{T}.
\]
Rosenstein selecciona para cada \(Y_i\) un vecino temporalmente separado
\(Y_{j(i)}\) y calcula \cite{Rosenstein1993}
\[
 d(k)=\left\langle
 \log\|Y_{i+k}-Y_{j(i)+k}\|_2
 \right\rangle_i,\qquad
 \widehat\lambda_{\max}=
 \frac1{\Delta t}\frac{\mathrm d d(k)}{\mathrm d k}.
\]
Eckmann ajusta mapas lineales locales de la forma
\[
 Y_{i+1}-Y_{j+1}\approx
 A_i(Y_i-Y_j)
\]
y obtiene un espectro finito mediante productos y
reortogonalización \cite{Eckmann1986}. Ordenados
\(\lambda_1\geq\cdots\geq\lambda_m\), la dimensión de Kaplan--Yorke es
\cite{Frederickson1983}
\[
 D_{\rm KY}=j+
 \frac{\sum_{i=1}^{j}\lambda_i}{|\lambda_{j+1}|},
 \qquad
 j=\max\left\{r:\sum_{i=1}^{r}\lambda_i\geq0\right\}.
\]
La función conserva intervalo de muestreo, unidades, parámetros de
reconstrucción, versión del backend, calidad del ajuste, advertencias y un
límite explícito para la matriz de distancias. La reconstrucción depende del
observable, del muestreo, de \(m\), \(\tau\), la exclusión temporal y la
ventana de ajuste.

\begin{verbatim}
from hidden_attractors import estimate_time_series_lyapunov

result = estimate_time_series_lyapunov(
    signal,
    sample_interval=0.01,
    time_unit="s",
    observable="x",
    rosenstein_emb_dim=8,
    eckmann_emb_dim=9,
    eckmann_matrix_dim=3,
)
print(result.largest_exponent, result.exponent_unit)
print(result.spectrum, result.kaplan_yorke_dimension)
\end{verbatim}

\subsection{Medidas de complejidad mediante backends opcionales}

\codepath{compute_complexity_measures} es un adaptador: no reimplementa
\codepath{nolds} ni \codepath{antropy}. Según el backend instalado puede
solicitar:
\begin{itemize}
  \item entropía muestral, \(\operatorname{SampEn}(m,r)
  =-\log(A_{m+1}/B_m)\) \cite{RichmanMoorman2000};
  \item dimensión de correlación,
  \(D_2=\lim_{r\to0}\mathrm d\log C(r)/\mathrm d\log r\)
  \cite{GrassbergerProcaccia1983};
  \item exponente de Rosenstein y exponente de Hurst
  \cite{Hurst1951};
  \item DFA, \(F(s)\propto s^\alpha\) \cite{Peng1994};
  \item entropía de permutación normalizada \cite{BandtPompe2002} y
  dimensión fractal de Higuchi \cite{Higuchi1988}.
\end{itemize}
Las convenciones exactas, valores predeterminados y correcciones pertenecen a
la versión registrada del backend. Por ello el resultado debe citar tanto la
librería adaptadora como el backend que realizó el cálculo.

\begin{verbatim}
from hidden_attractors.integrations import (
    available_complexity_backends,
    compute_complexity_measures,
)

print(available_complexity_backends())
complexity = compute_complexity_measures(
    signal,
    backend="auto",
    sample_rate=100.0,
    measures=[
        "sample_entropy",
        "permutation_entropy",
        "spectral_entropy",
        "dfa",
    ],
)
\end{verbatim}

\subsection{ADM Wu2023: ecuación implementada y frontera científica}

La ruta \codepath{adm_wu2023_integrate} es una reproducción específica del
Chua arctan de Wu et al. \cite{Wu2023ArctanChua}. En cada paso aproxima
\[
 X_{n+1}=X_n+
 \sum_{r=1}^{4}C^{(r)}
 \frac{h^{rq}}{\Gamma(rq+1)}.
\]
Sea \(g(x)=\arctan(\rho x)\), \(s=1+\rho^2x_n^2\), y
\[
 g'_n=\frac{\rho}{s},\quad
 g''_n=-\frac{2\rho^3x_n}{s^2},\quad
 g'''_n=\frac{8\rho^5x_n^2}{s^3}-\frac{2\rho^3}{s^2}.
\]
El primer coeficiente es \(C^{(1)}=F(X_n)\). Para la primera componente, los
polinomios de Adomian usados por el código son
\[
\begin{aligned}
 A_1={}&g'_n C^{(1)}_1,\\
 A_2={}&g'_n C^{(2)}_1+
 \frac{g''_n}{2}(C^{(1)}_1)^2
 \frac{\Gamma(2q+1)}{\Gamma(q+1)^2},\\
 A_3={}&g'_n C^{(3)}_1+
 g''_nC^{(1)}_1C^{(2)}_1
 \frac{\Gamma(3q+1)}{\Gamma(q+1)\Gamma(2q+1)}
 +\frac{g'''_n}{6}(C^{(1)}_1)^3
 \frac{\Gamma(3q+1)}{\Gamma(q+1)^3}.
\end{aligned}
\]
Para \(r=1,2,3\), el siguiente nivel satisface
\[
\begin{aligned}
 C^{(r+1)}_1&=-\alpha(1+a_1)C^{(r)}_1
 +\alpha C^{(r)}_2-\alpha a_2 A_r,\\
 C^{(r+1)}_2&=C^{(r)}_1-C^{(r)}_2+C^{(r)}_3,\\
 C^{(r+1)}_3&=-\beta C^{(r)}_2-\gamma C^{(r)}_3.
\end{aligned}
\]
Cada actualización usa únicamente \(X_n\). No contiene la convolución de
historia de Caputo y no es equivalente a ABM o EFORK-3 con memoria completa.
Su uso defendible es la reproducción metodológica, no sustituir una
integración Caputo.

\begin{verbatim}
import numpy as np
from hidden_attractors.integrations import adm_wu2023_integrate

times, states, status, info = adm_wu2023_integrate(
    params={
        "alpha": 10.0, "beta": 14.87, "gamma": 0.0,
        "a1": -1.27, "a2": 1.08, "rho": 1.0,
    },
    x0=np.array([0.1, 0.0, 0.0]),
    q=0.99,
    h=0.005,
    N=2000,
)
\end{verbatim}

\subsection{Integradores: selector y llamadas}

\begin{center}
\footnotesize
\begin{longtable}{@{}L{0.25\textwidth}L{0.16\textwidth}L{0.25\textwidth}L{0.26\textwidth}@{}}
\toprule
Entrada & Orden & Llamada mínima & Contrato \\
\midrule
\codepath{integrate(..., integrator="rk4")} & \(q=1\) &
\codepath{integrate(rhs,x0,1,h,tf,"rk4")} & RK4 clásico.\\
\codepath{integrate(..., integrator="heun")} & \(q=1\) &
\codepath{integrate(rhs,x0,1,h,tf,"heun")} & Trapecio explícito de orden dos.\\
\codepath{integrate(..., integrator="efork_q1")} & \(q=1\) &
\codepath{integrate(rhs,x0,1,h,tf,"efork_q1")} & Límite entero de coeficientes EFORK.\\
\codepath{integrate(..., integrator="abm")} & \(0<q<1\) &
\codepath{integrate(rhs,x0,q,h,tf,"abm")} & ABM Caputo; declarar memoria.\\
\codepath{integrate(..., integrator="efork3")} & \(0<q\leq1\) &
\codepath{integrate(rhs,x0,q,h,tf,"efork3")} & EFORK-3; en \(q=1\) redirige al límite entero.\\
\codepath{adm_wu2023_integrate} & \(0<q\leq1\) &
\codepath{adm_wu2023_integrate(params,x0,q,h,N)} & Reproducción arctan local, sin memoria Caputo.\\
\codepath{integrate_fractional_abm} & órdenes por componente &
\codepath{integrate_fractional_abm(rhs,x0,orders,h,N)} & Ruta experimental de reinicio por bloques; no es contrato público general no conmensurado.\\
\bottomrule
\end{longtable}
\end{center}

Las funciones de bajo nivel \codepath{integrate_general},
\codepath{caputo_abm_integrate}, \codepath{efork_integrate},
\codepath{rk4_integrate} y \codepath{fractional_integrate} existen para control
avanzado y pruebas. Para código nuevo se prefiere \codepath{integrate}, salvo
la reproducción ADM, que requiere sus parámetros específicos y se invoca de
forma directa.
